\section{What changes when the tensor reaches infinity}
\label{b:transition}

The boundary construction controls two different interfaces of the same
deformation.  At the inner boundary it prescribes the normal flux so that
the new mean curvature is negative.  At infinity it estimates the total
flux so that the new energy is no larger than the old energy, up to an
error tending to zero.  Compact support of the second fundamental form
simplifies the second task: sufficiently far out, the trace equation
contains only the graph.  We now explain precisely how maximality replaces
that simplification.  This calculation is also the reason to retain a
separate maximal construction: it carries the original decaying tensor
through the deformation instead of removing it from the end.

We work in spatial dimension four.  Write $g$ for the background metric,
$K$ for its second fundamental form, and assume $\tr_gK=0$.
Let $N\ge4$ be a fixed deformation parameter, let $0<\varepsilon<1$,
and choose a smooth nonincreasing function $\vartheta$ equal to one on
$(-\infty,0]$ and zero on $[1,\infty)$.  Set
\[
 p(t)=N\vartheta(Nt),\qquad
 l(t)=\exp\!\left(\int_0^t p(s)\,\dd s\right),\qquad
 L_0(t)=e^{2t}l(t),\qquad
 \ell=e^{-\varepsilon N},\quad \eta_N=\ell^{3/2}.
\]
For functions $f,Z$, the equation
\[
 Z=t+\tfrac16\log D,\qquad D=1+l(t)^2|\nabla f|_g^2
\]
defines $t$ uniquely: its right side has positive $t$ derivative
$1+p(1-D^{-1})/3$, and tends to the two ends of the real line with $t$.
Put
\[
 d=D^{-1},\quad w=\frac{l\nabla f}{\sqrt D},\quad v=|w|^2=1-d,
 \quad\chi=\frac{3d}{3+pv},\quad
 A_\chi=I-\frac{3+p}{3+pv}w\otimes w,
 \quad u=L_0\sqrt D.
\]
The deformed metric and flux are
\[
 \widehat g=e^{2t}(g+l^2\dd f^2),\qquad
 V=uA_\chi\bigl(3\nabla Z+K(w,\cdot)\bigr).
\]
These are exactly the four-dimensional boundary variables.  Thus the
inner boundary calculation remains unchanged.  On a face where $f$ is
constant, writing $H=\operatorname{div}_B\nu$ and $P_B=\tr_BK$ for
the normal into the exterior gives
\begin{equation}\label{b:transition:boundary}
 \frac{V_\nu}{u}
 =d(H+3\partial_\nu t)-H+w_\nu(F-P_B),
 \qquad
 \widehat H=e^{-t}\sqrt d\,(H+3\partial_\nu t).
\end{equation}
These are the local identities of
Lemma~\ref{four:sys:boundary-lemma}, with
$H^f=lD^{-1/2}\nabla^2f$ and $F=\tr_{A_\chi}(K+H^f)$.
All graph variables and the normal convention agree with that lemma.
Its proof uses only the constant trace of $f$ on the face and the
four-dimensional conformal boundary formula; it imposes no support
or energy hypothesis on $K$.
In particular the prescribed flux
$V_\nu/u=-H+w_\nu(F-P_B)-N_Bd$, with $N_B>0$, still makes
$\widehat H=-e^{-t}\sqrt d\,N_B<0$.

The new issue lies at the other end of the domain.  There the compactly
supported offset $C$ vanishes and the trace equation is
\begin{equation}\label{b:transition:trace}
 \tr_{A_\chi}(K+H^f)=\eta_N f.
\end{equation}
Maximality is useful here for an exact algebraic reason:
\[
 \tr_{A_\chi}K
 =-\frac{3+p}{3+pv}\frac{l^2}{D}K(\nabla f,\nabla f).
\]
Consequently \eqref{b:transition:trace} is, along any fixed solution,
the homogeneous linear equation
\begin{equation}\label{b:transition:drift}
 A_\chi:\nabla^2 f+b\cdot\nabla f-c_f f=0,
 \qquad
 b=-\frac{3+p}{3+pv}\frac l{\sqrt D}K(\nabla f,\cdot),
 \quad c_f=\frac{\eta_N\sqrt D}{l}\ge\frac{\eta_N}{e}>0.
\end{equation}
Here $l\le e$ follows directly from its definition.  Thus retaining $K$
introduces a drift, but no independent forcing of the graph.  An arbitrary
nonmaximal tensor would leave the term $\tr_gK$ and would not have this
property.

\begin{lemma}[The maximal end equation]\label{b:transition:end}
Suppose $1<q<2$, $0<\delta<q-1$,
$g-\delta_{\mathrm{Eucl}}=O_2(r^{-q})$ and
$K=O_1(r^{-1-q})$ on a four-dimensional end.  Assume
$\tr_gK=0$ and $R_g=O(r^{-4-\delta})$.
For fixed $N,\varepsilon$, consider smooth solutions on successively
larger spherical truncations, with $f=Z=0$ on the outer sphere, of
\eqref{b:transition:trace} and
\begin{equation}\label{b:transition:divergence}
 \operatorname{div}_gV
 =u\left[\delta_0\mathcal T+\rho+
       \delta_0\eta_N |w||\nabla f|
       -m_0(t)C_N(\rho_0+v\rho_1)\right].
\end{equation}
Here $0<\delta_0<1$, $\rho,\rho_0=O(r^{-4-\delta})$ with
their differentiated bounds, $\rho_1=O(r^{-2\gamma})$ for a fixed
$\gamma>1$, $m_0$ is a smooth bounded function, and $\mathcal T$ is
the quadratic term in the four-dimensional scalar identity.
Assume that the solutions have a common bounded range and that, for every
fixed derivative order, the background coefficients, weights and
solutions have uniform fixed-$N$ classical bounds on end unit patches,
including the outer boundary patches.  Assume also that $A_\chi$ is
uniformly elliptic there.  These bounds may depend on $N$ and the
derivative order; they are independent of the outer radius.  They are
separate assumptions from the displayed $O_2/O_1$ decay rates.
Then, uniformly on those truncations,
\[
 |\nabla^j f|\le C_j(N)e^{-c_j(N)r}\quad(j\ge0),
 \qquad |Z|+|t|\le C(N)r^{-2}.
\]
Every smooth local limit satisfies $Z,t=O_2(r^{-2})$.  Before taking
that limit, its scalar equation is
\begin{equation}\label{b:transition:zero-graph}
 \begin{split}
 3\Delta_g Z
 +3[p(Z)+2-\delta_0(p(Z)+1)]|\dd Z|^2
 ={}&\tfrac{\delta_0}{2}|K|^2+\rho
       -m_0(Z)C_N\rho_0+\mathcal E_N,
 \end{split}
\end{equation}
where $\mathcal E_N$ and each required fixed derivative have
exponential unit-patch bounds.
\end{lemma}

\begin{proof}
The coefficients in \eqref{b:transition:drift} are uniformly bounded
and elliptic at fixed $N$.  For sufficiently small $c>0$, substitution
of $e^{-cr}$ into its left side gives an upper bound
\[
 e^{-cr}\bigl(C(N)c^2+C(N)c-\eta_N/e\bigr)<0
\]
outside one fixed sphere.  A multiple of this positive supersolution
dominates $|f|$ on that sphere; the zero outer value is already dominated.
Comparison for $f$ and $-f$ gives exponential decay, uniformly in the
truncation radius.  The assumed unit-patch estimates, followed by the
smooth linear interior and zero-Dirichlet boundary estimates for
\eqref{b:transition:drift}, give the asserted differentiated decay.
These are the smooth-coefficient estimates of
\cite[Chs.~6--9]{GilbargTrudinger2001} and
\cite{AgmonDouglisNirenberg1959}.
The implicit formula for $t$ then shows that $t-Z$ and the graph metric
error have the same type of bounds.

It remains to identify what survives when those graph errors are removed.
Set the graph height and its first two derivatives equal to zero in
the coefficients, retaining the exponentially small difference as an
error.  Then $H^f=0$, $F=0$, $w=0$, $D=d=\chi=1$, $t=Z$ and
$u=L_0(Z)$.  The smooth scalar quadratic form at this zero graph is
\begin{equation}\label{b:transition:quadratic}
 \mathcal T=\tfrac12|K|^2+3(p(Z)+1)|\dd Z|^2.
\end{equation}
One can also verify this directly from its block formula.  Choose any
unit vector $e$, write $j=K(e,e)$, $M=K(e,\cdot)|_{e^\perp}$ and
$T=K|_{e^\perp}$, and use $\tr T=-j$.  The terms involving $K$ become
$\tfrac12|T^0|^2+|M|^2+\tfrac23j^2$.
Since $|T|^2=|T^0|^2+j^2/3$, this is exactly $|K|^2/2$.
The terms linear in $\dd Z$ contain $w$ and vanish, leaving the
second term in \eqref{b:transition:quadratic}.  In particular setting
the graph to zero does not set the tensor to zero.

Now $(\log L_0)'=p+2$.  Divide
\eqref{b:transition:divergence} by $L_0$, insert
\eqref{b:transition:quadratic}, and move the quadratic gradient term
to the left.  Smooth dependence on the bounded solution variables
absorbs the graph errors into $\mathcal E_N$ and gives
\eqref{b:transition:zero-graph}.

Define $\Phi(0)=0$ and
\[
 \Phi'(z)=\exp\!\left(\int_0^z
     [p(s)+2-\delta_0(p(s)+1)]\,\dd s\right).
\]
On the common bounded range of $Z$, both $\Phi'$ and its reciprocal
are bounded at fixed $N$.  The chain rule gives
\[
 \Delta_g\Phi(Z)=\frac{\Phi'(Z)}3
 \left[\tfrac{\delta_0}{2}|K|^2+\rho
       -m_0(Z)C_N\rho_0+\mathcal E_N\right].
\]
The right side is $O_N(r^{-4-\delta})$, because
$|K|^2=O(r^{-2-2q})$ and $2+2q>4+\delta$.
With $\alpha=\delta/2$, set $b_0=r^{-2}-r^{-2-\alpha}$.
For sufficiently large $r$, this is positive and
\[
 -\Delta_g b_0
 =\alpha(2+\alpha)r^{-4-\alpha}+O(r^{-4-q})
 \ge c r^{-4-\alpha}.
\]
A fixed multiple of $b_0$ dominates both the absolute source and
$|\Phi(Z)|$ on the fixed inner sphere.  Since $\Phi(Z)=0$ at the outer
sphere, comparison for both signs gives $|\Phi(Z)|\le C(N)r^{-2}$,
and hence $|Z|+|t|\le C(N)r^{-2}$.  These estimates hold before
exhaustion, so the local limit has the asserted normalization at
infinity.

Finally rescale annuli $R<r<2R$ to unit size.  The metric has uniformly
controlled scaled $C^{1,\beta}$ coefficients for each fixed $0<\beta<1$.
Interior $W^{2,p_1}$ estimates with $p_1>4$ first give
$|\dd Z|=O_N(r^{-3})$.  The first derivative bound on $K$, the
differentiated weights and the exponential error then give a controlled
scaled H\"older norm for the source.  Schauder estimates yield
$\nabla^2Z=O_N(r^{-4})$.  The same follows for $t$ from the graph error.
\end{proof}

The term $|K|^2/2$ has two consequences for the rest of the construction.
First, the end proof must use the original decay of $K$; it cannot invoke
a result whose hypothesis is $K=0$ outside a compact set.  Second, the
continuation must control the corresponding noncompact drift before
these fixed-parameter estimates are available.  The appropriate tail
weights follow from a short calculation.  Choose
\[
 1<\gamma<\min\{2,q-\delta/2\},\qquad
 \rho_0\asymp r^{-4-\delta},\quad \rho_1\asymp r^{-2\gamma}.
\]
This interval is nonempty because $\delta<q-1$.
For $a=|w|$ one has
\begin{equation}\label{b:transition:tail-budget}
 a r^{-2-q}
 \le\tfrac12a^2r^{-2\gamma}
       +\tfrac12r^{-4-2q+2\gamma}
 \le C(v\rho_1+\rho_0).
\end{equation}
The first inequality is Young's inequality; the second uses
$2\gamma<2q-\delta$.  Also $|K|^2\le C\rho_0$.
Thus the drift error controlled by
$\Pi_N(|K|^2+a|\nabla K|)$ fits the same kind of floor penalty as
in the boundary construction, with exponents adapted to the original
end.  Since $4\gamma>4$, the square of $\rho_1$ is integrable.
That last fact is exactly what permits its contribution to the flux
loss to be absorbed with an error tending to zero.

The local comparison is now explicit.  Signed boundary faces and their
inner flux formula are shared.  The maximal route supplies its own
strict approximation and an all-black threshold boundary; it proves
the noncompact drift estimate and its complete boundary continuation
with the weights above.  Once these yield the fixed-$N$ estimates,
Lemma~\ref{b:transition:end} gives the correct end normalization.
In applying that lemma, unit-scale convolution of the strict end
coefficients supplies their uniform bounds at every fixed derivative
order; a subsequent trace-free projection preserves maximality.  This
is an actual preparation step, not a consequence of smoothness and
$O_2/O_1$ decay alone.  The displayed decay rates are the ones retained
from the original data and used in the final annular estimates.
The final comparison still uses the actual metric $\widehat g$, its
minimal enclosure, and its ADM flux.  Maximality preserves a method
with the original tensor tail; it does not replace those existence,
boundary, or enclosure arguments.
